% Shared preamble for the English summaries in this repository.
\documentclass[11pt,a4paper]{article}
\usepackage[margin=2.2cm]{geometry}
\usepackage[T1]{fontenc}
\usepackage[utf8]{inputenc}
\usepackage{lmodern}
\usepackage{amsmath,amssymb}
\usepackage{graphicx}
\usepackage{booktabs}
\usepackage{enumitem}
\usepackage[hidelinks]{hyperref}
\setlist{itemsep=2pt,topsep=4pt}
\setlength{\parindent}{0pt}
\setlength{\parskip}{5pt}
\graphicspath{{figures/}}
\newcommand{\summaryheader}[3]{%
  \begin{center}
    {\Large\bfseries #2}\\[3pt]
    {\large #1}\\[3pt]
    {\small Arshia Goshtasbi \textbullet{} MSc Electrical Engineering (Control), Iran University of Science and Technology \textbullet{} #3}
  \end{center}
  \vspace{2pt}\hrule\vspace{8pt}
  \noindent\textit{English summary of my original Persian report, which is kept as a PDF in this folder.}\par}

\begin{document}
\summaryheader{Optimal Control, Homework 3}{Constrained Extrema and Lagrange Multipliers}{Spring 1404 (2025)}

\noindent\textit{The problems come from Chapter~4 of Kirk, \emph{Optimal Control Theory}.}

\begin{description}
  \item[Problem 4-22 (static constrained extremum).] I extremised $f=y_1^2+y_2^2+y_3^2$ subject to $y_1+y_2+y_3=5$ and $y_1^2+y_2^2+y_3=9$. I formed the augmented function with two multipliers and set its differential to zero, then examined the case where $y_1^\star=y_2^\star$ and $p_2^\star=-1$. This gives the candidates $(2,2,1)$ with $f=9$ and $(-1,-1,7)$ with $f=51$. The other case ($p_1^\star=0$) has a negative discriminant and no real solution.
  \item[Problem 4-24 (differential constraints).] For $J=\int(\lambda+w_3^2)\,dt$ with $\dot w_1=w_2$ and $\dot w_2=w_3$, I wrote the augmented integrand $g_a=\lambda+w_3^2+p_1(\dot w_1-w_2)+p_2(\dot w_2-w_3)$. The three Euler--Lagrange equations give $p_1^\star$ constant, $\dot p_2^\star=-p_1^\star$ and $p_2^\star=2w_3^\star$.
  \item[Problem 4-25 (isoperimetric constraint).] To minimise $\int_0^1(\dot x^2+t^2)\,dt$ with $x(0)=x(1)=0$ and $\int_0^1 x^2\,dt=2$, I introduced $z$ with $\dot z=x^2$. The Euler--Lagrange equations give $\ddot x=Cx$ with constant $C$. Only $C=-n^2\pi^2$ admits a non-trivial solution satisfying both end conditions, and the integral constraint fixes the amplitude: $x^\star(t)=2\sin(n\pi t)$.
\end{description}
\end{document}
